% titel: Blatt 7, A1: Lewis-Formeln mit Oktettregel und erweiterter Valenzschale
% skript: Mortimer Kap. 9 und 10 (offene Oktette)

% Teil I: Werkzeuge, kein L\"osungsweg.

Ziel ist es, f\"ur $\mathrm{POCl_3}$ und $\mathrm{H_2SO_4}$ je zwei
vollst\"andige Lewis-Formeln zu zeichnen: eine, in der jedes Atom die
Oktettregel erf\"ullt, und eine, in der das Zentralatom die Valenzschale
erweitert (Hypervalenz/offenes Oktett).

\begin{werkzeugkasten}[Vorgehen f\"ur Lewis-Formeln]
\begin{itemize}
  \item Zahl der Valenzelektronen bestimmen (Gruppennummer $+$ Ladung bei
        Ionen), durch 2 teilen $\rightarrow$ Anzahl Elektronenpaare.
  \item Zentralatom ist das am wenigsten elektronegative Nichtmetall.
        $\mathrm{H}$ und Halogene kommen als Endatome vor, nicht als
        Zentralatom.
  \item Einfachbindungen entwerfen, freie Elektronenpaare auf das Oktett
        verteilen (erst Endatome, dann zentral).
  \item Formal\-ladungen pr\"ufen: Formalladung = Valenzelektronen -
        (freie Paare) - 0{,}5$\times$(Bindungselektronen).
  \item Falls ein Zentralatom der 3. Periode oder h\"oher (z.~B. P, S) zu
        wenig Oktett-Partner hat, k\"onnen Doppelbindungen oder erweiterte
        Valenzschalen (d-Orbitale) hinein.
\item \textbf{Oktett-Form:} alle Atome mit 8 Valenzelektronen, evtl. mit
      Formalladungen, keine Doppelbindungen zum Zentralatom n\"otig.
  \item \textbf{Erweiterte Valenzschale:} Zentralatom (P, S) bindet mehr
        als 4 Bindungspaare, oft mehrere Doppelbindungen $\rightarrow$
        Formalladungen oft 0.
\end{itemize}
\textit{Skript:} Mortimer Kap. 9 (Lewis), Kap. 10 (Oktett und offene
Oktette).
\end{werkzeugkasten}

\begin{achtung}[Darauf achten]
\begin{itemize}
  \item P (Gruppe 15) und S (Gruppe 16) liegen in der 3. Periode; sie
        k\"onnen mehr als ein Oktett an Bindungspaaren aufnehmen (d-Orbitale).
  \item In der Oktett-erf\"ullenden Form von $\mathrm{POCl_3}$ tragen P und
        O Formalladungen; die erweiterte Form (P=O-Doppelbindung) hat alle
        Formalladungen 0.
  \item Bei $\mathrm{H_2SO_4}$ (S umgeben von 4 O) erzwingt die Oktettregel
        zwei Doppelbindungen O=S=O; h\"ohere S-Bindungszahl $=$ erweiterte
        Valenzschale.
\end{itemize}
\end{achtung}
